\inicializarSeccion{Funcionalidad de las palabras claves \texttt{\textbar for\textbar} y \texttt{\textbar while\textbar}}
\tab Los ciclos for y while son estructuras de control que permiten repetir un bloque de instrucciones de código varias veces sin necesidad de escribirlo repetidamente. Tienen la principal función de automatizar procesos iterativos.

Se utilizan cuando:
\begin{itemize}
    \item Se necesita realizar cálculos repetitivos (por ejemplo, sumar contribuciones de varias cargas eléctricas).
    \item Se recorren arreglos o listas de datos.
    \item Se ejecuta un proceso hasta cumplir cierta condición.
\end{itemize}

En el contexto del cálculo del campo eléctrico, estos ciclos son muy importantes para sumar las contribuciones de múltiples cargas, ya que cada carga aporta un campo que se debe calcular individualmente y luego sumarse.